\documentclass[12pt]{article}
\usepackage{fontspec}
\usepackage[english, russian]{babel}
\usepackage{amsmath, amssymb}
\usepackage{tikz}
\usepackage{pgfplots}
\usepackage{longtable}
\usepackage{lscape}
\pgfplotsset{compat=1.18}

\setmainfont{CMU Serif}
\setsansfont{CMU Sans Serif}
\setmonofont{CMU Typewriter Text}

\usepackage[margin=2cm]{geometry}
\begin{document}
\section{Метод квадратичной аппроксимации}
\subsection{Применимость метода}
Дана функция одной переменной \( f(x) = x^2-3x+x\ln x \), определённая на отрезке \( [0.9, 2.0] \).
Метод применим, так как функция предполагается унимодальной и непрерывной на заданном отрезке, что гарантирует существование единственного минимума, сходимость алгоритма и возможность аппроксимации многочленом.
\subsection{График исходной функции и реальный минимум}
Точка минимума из графика: \( x_* \approx 0.999968 \), \( f(x_*) \approx -2 \).
\begin{center}
\begin{tikzpicture}
\begin{axis}[
  width=12cm, height=8cm,
  grid=major,
  grid style={dashed, black!30},
  legend pos=outer north east,
  enlargelimits=true,
  xlabel={\(x\)},
  ylabel={\(f(x)\)}
]
\addplot [domain=0.9:1.45, samples=100, dashed, thick] {x^2 - 3*x + x*ln(x)};
\addlegendentry{Исходная функция}
\addplot [domain=0.6799999999999999:2.22, samples=100, solid, thick] {1.3536125345018937*x^2 + -2.679004690771915*x + -0.6701463953438525};
\addlegendentry{Аппроксимация}
\addplot[only marks, mark=*, text=black] coordinates {(0.9, -1.9848244640920438) (1.45, -1.7087328431728994) (2.0, -0.6137056388801094)};
\addlegendentry{Узлы}
\end{axis}
\end{tikzpicture}
\end{center}
\subsection{Начальные значения}
Выберем начальные точки так, чтобы выполнялось условие \( f_2 \le f_1 \) и \( f_2 \le f_3 \):
\begin{itemize}
  \item \( x_1^{(1)} = 0.9, \; f_1^{(1)} = -1.9848245 \)
  \item \( x_2^{(1)} = 1.10108, \; f_2^{(1)} = -1.9848382 \)
  \item \( x_3^{(1)} = 2, \; f_3^{(1)} = -0.6137056 \)
\end{itemize}
\subsection{Итерация 1}
Проверка критерия останова: интервал неопределенности равен \( 2 - 0.9 = 1.1 \).
Так как \( 1.1 \ge 0.0001 \), продолжаем вычисления.


Вычисляем \( \tilde{x}_{*}^{(1)} \):
\begin{multline*}
\tilde{x}_{*}^{(1)} = \frac{1}{2}\cdot\frac{f_1(x_2^2 - x_3^2) + f_2(x_3^2 - x_1^2) + f_3(x_1^2 - x_2^2)}{f_1(x_2 - x_3) + f_2(x_3 - x_1) + f_3(x_1 - x_2)} = \\
= \frac{1}{2}\cdot\frac{-1.9848245(-2.7876228) + -1.9848382(3.19) + -0.6137056(-0.4023772)}{-1.9848245(-0.89892) + -1.9848382(1.1) + -0.6137056(-0.20108)} = 1.0005646
\end{multline*}
Вычисляем значение функции: \( \tilde{f}_{*}^{(1)} = f(1.0005646) = -1.9999995\).
Проверяем условия выбора новой тройки точек:
\begin{itemize}
  \item а) \( \tilde{x}_{*}^{(1)} \in [x_2, x_3] \) и \( \tilde{f}_{*}^{(1)} \le f_2^{(1)} \) --- Ложно.
  \item б) \( \tilde{x}_{*}^{(1)} \in [x_2, x_3] \) и \( \tilde{f}_{*}^{(1)} > f_2^{(1)} \) --- Ложно.
  \item в) \( \tilde{x}_{*}^{(1)} \in [x_1, x_2] \) и \( \tilde{f}_{*}^{(1)} \le f_2^{(1)} \) --- Истинно.
  \item г) \( \tilde{x}_{*}^{(1)} \in [x_1, x_2] \) и \( \tilde{f}_{*}^{(1)} > f_2^{(1)} \) --- Ложно.
\end{itemize}
Выполнено условие в).
\begin{center}
\begin{tikzpicture}
\begin{axis}[
  width=12cm, height=8cm,
  grid=major,
  grid style={dashed, black!30},
  legend pos=outer north east,
  enlargelimits=true,
  xlabel={\(x\)},
  ylabel={\(f(x)\)}
]
\addplot [domain=0.9:1.10108, samples=100, dashed, thick] {x^2 - 3*x + x*ln(x)};
\addlegendentry{Исходная функция}
\addplot [domain=0.6799999999999999:2.22, samples=100, solid, thick] {1.386708299729599*x^2 + -2.7749824099322593*x + -0.6105740179339832};
\addlegendentry{Аппроксимация}
\addplot[only marks, mark=*, text=black] coordinates {(0.9, -1.9848244640920438) (1.10108, -1.9848381708126648) (2.0, -0.6137056388801094)};
\addlegendentry{Узлы}
\addplot[only marks, mark=square*, text=black] coordinates {(1.000564578171691, -1.9999995219072169)};
\addlegendentry{\(\tilde{x}_{*}^{(1)}\)}
\end{axis}
\end{tikzpicture}
\end{center}
\subsection{Итерация 2}
Проверка критерия останова: интервал неопределенности равен \( 1.10108 - 0.9 = 0.20108 \).
Так как \( 0.20108 \ge 0.0001 \), продолжаем вычисления.


Вычисляем \( \tilde{x}_{*}^{(2)} \):
$$ \tilde{x}_{*}^{(2)} = \frac{1}{2}\cdot\frac{-1.9848245(-0.2112477) + -1.9999995(0.4023772) + -1.9848382(-0.1911295)}{-1.9848245(-0.1005154) + -1.9999995(0.20108) + -1.9848382(-0.1005646)} = 1.0005627 $$
Вычисляем значение функции: \( \tilde{f}_{*}^{(2)} = f(1.0005627) = -1.9999995\).
Проверяем условия выбора новой тройки точек:
\begin{itemize}
  \item а) \( \tilde{x}_{*}^{(2)} \in [x_2, x_3] \) и \( \tilde{f}_{*}^{(2)} \le f_2^{(2)} \) --- Ложно.
  \item б) \( \tilde{x}_{*}^{(2)} \in [x_2, x_3] \) и \( \tilde{f}_{*}^{(2)} > f_2^{(2)} \) --- Ложно.
  \item в) \( \tilde{x}_{*}^{(2)} \in [x_1, x_2] \) и \( \tilde{f}_{*}^{(2)} \le f_2^{(2)} \) --- Истинно.
  \item г) \( \tilde{x}_{*}^{(2)} \in [x_1, x_2] \) и \( \tilde{f}_{*}^{(2)} > f_2^{(2)} \) --- Ложно.
\end{itemize}
Выполнено условие в).
\begin{center}
\begin{tikzpicture}
\begin{axis}[
  width=12cm, height=8cm,
  grid=major,
  grid style={dashed, black!30},
  legend pos=outer north east,
  enlargelimits=true,
  xlabel={\(x\)},
  ylabel={\(f(x)\)}
]
\addplot [domain=0.9:1.000564578171691, samples=100, dashed, thick] {x^2 - 3*x + x*ln(x)};
\addlegendentry{Исходная функция}
\addplot [domain=0.859784:1.141296, samples=100, solid, thick] {1.5005704510892155*x^2 + -3.002829683774963*x + -0.4977398140768306};
\addlegendentry{Аппроксимация}
\addplot[only marks, mark=*, text=black] coordinates {(0.9, -1.9848244640920438) (1.000564578171691, -1.9999995219072169) (1.10108, -1.9848381708126648)};
\addlegendentry{Узлы}
\addplot[only marks, mark=square*, text=black] coordinates {(1.0005627131986043, -1.9999995250604723)};
\addlegendentry{\(\tilde{x}_{*}^{(2)}\)}
\end{axis}
\end{tikzpicture}
\end{center}
\subsection{Итерация 3}
Проверка критерия останова: интервал неопределенности равен \( 1.0005646 - 0.9 = 0.1005646 \).
Так как \( 0.1005646 \ge 0.0001 \), продолжаем вычисления.


Вычисляем \( \tilde{x}_{*}^{(3)} \):
$$ \tilde{x}_{*}^{(3)} = \frac{1}{2}\cdot\frac{-1.9848245(-0.0000037) + -1.9999995(0.1911295) + -1.9999995(-0.1911257)}{-1.9848245(-0.0000019) + -1.9999995(0.1005646) + -1.9999995(-0.1005627)} = 1.0000065 $$
Вычисляем значение функции: \( \tilde{f}_{*}^{(3)} = f(1.0000065) = -2\).
Проверяем условия выбора новой тройки точек:
\begin{itemize}
  \item а) \( \tilde{x}_{*}^{(3)} \in [x_2, x_3] \) и \( \tilde{f}_{*}^{(3)} \le f_2^{(3)} \) --- Ложно.
  \item б) \( \tilde{x}_{*}^{(3)} \in [x_2, x_3] \) и \( \tilde{f}_{*}^{(3)} > f_2^{(3)} \) --- Ложно.
  \item в) \( \tilde{x}_{*}^{(3)} \in [x_1, x_2] \) и \( \tilde{f}_{*}^{(3)} \le f_2^{(3)} \) --- Истинно.
  \item г) \( \tilde{x}_{*}^{(3)} \in [x_1, x_2] \) и \( \tilde{f}_{*}^{(3)} > f_2^{(3)} \) --- Ложно.
\end{itemize}
Выполнено условие в).
\begin{center}
\begin{tikzpicture}
\begin{axis}[
  width=12cm, height=8cm,
  grid=major,
  grid style={dashed, black!30},
  legend pos=outer north east,
  enlargelimits=true,
  xlabel={\(x\)},
  ylabel={\(f(x)\)}
]
\addplot [domain=0.9:1.0005627131986043, samples=100, dashed, thick] {x^2 - 3*x + x*ln(x)};
\addlegendentry{Исходная функция}
\addplot [domain=0.8798870843656619:1.0206774938060292, samples=100, solid, thick] {1.517355792370382*x^2 + -3.0347313087702714*x + -0.48262447945487547};
\addlegendentry{Аппроксимация}
\addplot[only marks, mark=*, text=black] coordinates {(0.9, -1.9848244640920438) (1.0005627131986043, -1.9999995250604723) (1.000564578171691, -1.9999995219072169)};
\addlegendentry{Узлы}
\addplot[only marks, mark=square*, text=black] coordinates {(1.0000064980496681, -1.9999999999366629)};
\addlegendentry{\(\tilde{x}_{*}^{(3)}\)}
\end{axis}
\end{tikzpicture}
\end{center}
\newpage
\begin{landscape}
\section{Таблица результатов}
\subsection{Остальные итерации}
\begin{longtable}{|c|c|c|c|c|c|c|c|}
\hline
\( k \) & \( x_1^{(k)} \) & \( x_2^{(k)} \) & \( x_3^{(k)} \) & \( \tilde{x}_{*}^{(k)} \) & \( f(\tilde{x}_{*}^{(k)}) \) & \( f_2^{(k)} \) & Условие \\
\hline
\endhead
4 & 0.9 & 1.00000649805 & 1.000562713199 & 1.000003291112 & -1.999999999984 & -1.999999999937 & в \\
\hline
5 & 0.9 & 1.000003291112 & 1.00000649805 & 1.000000056456 & -2 & -1.999999999984 & в \\
\hline
6 & 0.9 & 1.000000056456 & 1.000003291112 & 1.000000019224 & -2 & -2 & в \\
\hline
7 & 0.9 & 1.000000019224 & 1.000000056456 & 1 & -2 & -2 & в \\
\hline
8 & 0.9 & 1 & 1.000000019224 & 1 & -2 & -2 & а \\
\hline
9 & 1 & 1 & 1.000000019224 & 1 & -2 & -2 & а \\
\hline
\end{longtable}
\end{landscape}
\end{document}